%% kapa: exam=1 round=33 year=2022 period=1 form=A subject=civil_law subject_ko=민법 seq=21 number=21
%% source: Q-Net 감정평가사 기출문제 (cst00309) artlSeq=5209997 file=raw/1008/5209997/2022년도 제33회 감정평가사 1차시험 1교시(민법,경제학,부동산학).pdf
%% conversion: kordoc 4.24.51 -> tools/gen1.py (auto); review=proofread-sonnet; flags=none
\begin{question}{21}
부동산등기에 관한 설명으로 옳지 않은 것은? (다툼이 있으면 판례에 따름)
\begin{choices}
  \item 가등기된 권리의 이전등기는 가등기에 대한 부기등기의 형식으로 할 수 있다.
  \item 근저당권등기가 원인 없이 말소된 경우, 그 회복등기가 마쳐지기 전이라도 말소된 등기의 등기명의인은 적법한 권리자로 추정된다.
  \item 청구권보전을 위한 가등기에 기하여 본등기가 경료되면 본등기에 의한 물권변동의 효력은 가등기한 때로 소급하여 발생한다.
  \item 소유권이전등기의 원인으로 주장된 계약서가 진정하지 않은 것으로 증명되었다면 그 등기의 적법추정은 복멸된다.
  \item 동일 부동산에 관하여 등기명의인을 달리하여 중복된 소유권보존등기가 경료된 경우, 선행보존등기가 원인무효가 아닌 한 후행보존등기는 실체관계에 부합하더라도 무효이다.
\end{choices}
\end{question}
\answer{3}
